\section{An Interpretation of Genesis}

Many Christians assume a strongly ``literal'' stance upon the first chapter of Genesis. It should go without saying that assuming Genesis 1 to have a perspective like a newspaper is outside of the bounds of any known form of traditionalism, conceivable or not -- God did not dictate this to a female stenographer, nor did the Scriptures fall from the sky. What might Genesis 1 be attempting to actually convey to the reader (or, we should hope, to the initiate)?

For assistance in this, let us peruse a few quotes from Max Heindel\footnote{\url{https://en.wikipedia.org/wiki/Max\_Heindel}}, someone worth salvaging out of the wreck of various Theosophies. We have already seen how Mouravieff focuses attention upon the New Testament Gnosis, and how he tends to uncover the various truths behind the two separate races of original beings which populated the earth, as well as the two sexes which originally formed one being. As an aside, consider how inversion can be used by the occult -- just as modern man separates Time and unites Space\footnote{\url{http://www.gornahoor.net/?p=3166}}, as one would in opposite-world, so one might say that Mouravieff teaches that we have united the races and separated the sexes, the separation of the sexes here being taken in our own time in a metaphorical sense by pointing towards the detaching of women from their consorts .

By contrast with Mouravieff's New Testament (or Third Testament), Heindel attempts to elucidate the ``pre''-era of Genesis, the time of the archetypes and the archons. As a disclaimer, I will note that Heindel believes in reincarnation, but not necessarily in a form which is incapable of being tamed to the necessary strictures of traditionalism -- he accepts ``evolution'', but not in a purely Steinerian sense; that is, man is not always stuck in an esoteric trajectory that involves some kind of ``mass'' evolution, but can actually assert the power of his primordial being and potentially royal
\& imperial destiny, as an individual. Heindel may be a lens to salvage and save Steiner.

\begin{quotex}
It was during the months of April and early May 1908 that Max Heindel stood the test. Only later was he told that the Elder Brother's first choice for candidate was Dr. Rudolf Steiner, who had been under their instruction for several years but had failed to pass his test because he could not be a leader of both the Western Teaching as well as the Eastern Teaching. Max Heindel also had been under the observation of the Elder Brothers for a number of years as the next best candidate in case the first one should fail. In addition, he was told that the Teachings had to be published before the close of the first decade April 9, 1910 
\end{quotex}


Steiner\footnote{\url{http://www.rffriends.org/wpx/wp-content/uploads/2009/10/Microsoft-Word-Chapter-4-Heindel-in-Germany.pdf}} had failed some kind of test, and was mixing Western and Eastern esotericism.

\begin{quotex}
While the material which now forms the Earth was yet a part of the Sun, it was, of course, in a fiery condition; but as the fire does not burn spirit, our human evolution commenced at once, being confined particularly to the Polar region of the Sun. The highest evolved beings which were to become human were the first to appear.
\end{quotex}

Notice that it is implied that human beings are manifestations of various spirits, not all of whom are ``equal''.

\begin{quotex}
In the beginning of the Lemurian epoch, these ``failures'' (note that they were failures, not merely stragglers) had crystallized that part of the Earth occupied by them to such a degree that it become a huge cinder or clinker, in the otherwise soft and fiery Earth. They were a hindrance, and an obstruction, so they, with the part of the Earth they had crystallized, were thrown out into space beyond recall. That is the genesis of our Moon.
\end{quotex}

Not all humans are to be regarded as worthy of imitation. One recalls Bondarev's association of the moon with the epoch of time when Yahweh governed the tribes of Israel; Yahweh was a mask of the archangel Michael, a spirit of ``tutelage'' in the very Godhead itself.

\begin{quotex}
That such is the effect of the Sun is shown by the rapidity of growth at the tropics... on the other hand had the Moon remained with the Earth, man would have crystallized into a statue... the separation of the Earth from the Sun, which now sends its rays from a far distance, enables man to live at the proper rate of vibration, to unfold slowly... the Sun works in the vital body and is the force which makes for life, and wars against the death-dealing Moon force.
\end{quotex}

This is evolutionary cosmology integrated into a mythical background of Traditionalism.

\begin{quotex}
The Lords of Form vivified the human spirit in as many of the stragglers of the Moon Period as had made the necessary progress in the three and one half Revolutions which had elapsed since the commencement of the Earth Period, but at that time the Lords of Mind could not give them the germ of Mind. Thus a great part of nascent humanity was left without this link between the threefold spirit and the threefold body.
\end{quotex}

Here is Mouravieff's doctrine of the original pre-Adamic race which was anthropoid in relation to the sons of God and the demigods.

He goes on to assert that during the Hyperborean epoch, man was capable of generating another being from himself without physical intercourse, because he was bi-polar, sexually. Of course, he was also predominantly non-material.

\begin{quotex}
The beings who inhabit Venus and Mercury are not quite so far advanced as those whose present field of evolution is the Sun, but they are very much further advanced than our humanity. Therefore they stayed somewhat longer with the central mass than did the inhabitants of the Earth... some inhabitants of each planet were sent to Earth to help nascent humanity... the lords of Venus were leaders of the masses of our people. They were inferior beings of the Venus evolution, who appeared among men and were known as messengers of the Gods. For the good of our humanity they led and guided it, step by step. There was no rebellion against their authority, because man had not as yet evolved an independent will. Their commands were obeyed without question, and they were held in deep reverence. When under this tuition mankind had reached a certain stage... the most advanced were placed under the guidance of the Lords of Mercury, who initiated them into the higher truths for the purpose of making them leaders of the people. These initiates were then exalted to kingship and were the founders of the Divine dynasties of rulers who were indeed kings ``by the grace of God''. The present age cannot be called a golden age in any save a material sense, but it is perhaps a necessary one, in order to bring man to the point where he will be able to rule himself, for self-mastery is the end and all aim of all rulership. No man can safely remain ungoverned who has not learned to govern himself... at the present time, this is the hardest command that can be given, to govern the self. It is easy to govern others, it is hard to force obedience from one's self. The Lords of Mercury now work upon the individual(s).
\end{quotex}


One might perhaps note here that all provision is made both for the rise of monarchy, its inner and outer worth, as well as an explanation offered for its ``fall'', along with the only imaginable justification (which is that each man might learn to rule himself as a king swayed his peoples). This doesn't invalidate the possibility for monarchy even today, it merely offers an existential and cosmological perspective on why sacred kingship is so despised by most men, currently, and what the kernel of truth in that might be.

Heindel goes on to discuss the Lemurian epoch in this chapter of \emph{The Rosicrucian CosmoConception}, Chapter XII, Evolution on the Earth.

Hopefully, this exercise helps demonstrate some of the various continuities which persist among very ``different'' esoteric teachings, made possible when one discards a modernistic and literal approach to Sacred Scripture; Genesis 1 is teaching a symbolic truth in terms of myth, which pre-originates any physical reality that derives from it. It is also worth mentioning that Saint Paul explicitly teaches about the ``time of tutelage'' under ``spirits'' in Galatians, and indeed in other passages (e.g., Hebrews warns to be hospitable to strangers).


\flrightit{Posted on 2011-10-21 by Logres}

